\section{Measurement-Grounded Calibration: the S0--S3 Ladder}
\label{sec:calibration}

\subsection{Design principle: one owner per phenomenon}

The ladder adds one measured phenomenon at a time, and each level owns
exactly one. The rule matters because the obvious alternative --
injecting everything measured at every level -- double-counts: a
detection-probability model and a dropout-burst model both reduce the
number of observations that reach the planner, and applied together they
reduce it twice. We therefore fix that S1 owns the \emph{marginal}
probability of detection and S2 owns only the \emph{correlation
structure} of the misses, parameterised so as to reproduce S1's marginal
exactly. The safeguard is structural rather than a matter of care: the
two-state visible/blind gate is constructed so that its stationary
visible fraction equals the marginal by identity, so S2 changes when the
misses fall and never how many there are.

\subsection{S1: a static observation model}

S1 was fitted to paired samples in which the real detector's output and
an independently measured range to the anchor were recorded together
($n=24$ usable pairs). Three quantities were extracted. The
\emph{probability of detection} varies strongly with range and
non-monotonically: \num{0.67} below \SI{0.8}{\metre}, \num{0.94} between
\SI{1.6} and \SI{2.0}{\metre}, and \num{0.15} beyond \SI{2.0}{\metre}.
The \emph{monocular range estimator} reads \num{0.469} of true range with
a median absolute error of \SI{1.01}{\metre}; the bias is injected at its
physical origin, the apparent height, rather than as an additive error on
range. The \emph{bearing} error has a systematic component that is a
frame convention rather than a sensor property, so only its scatter
(\SI{0.32}{\degree}) is injected.

The dataset has a gap between \SI{0.8} and \SI{1.6}{\metre} that no
amount of care removes, and the detection-probability curve is therefore
interpolated across the band in which, as it turned out, the physical
runs spent most of their time.

\subsection{S2: rate and the structure of the gaps}

S2 was fitted to the timing of the same recorded stream. Observations
arrive at \SI{3.5}{\hertz}; the marginal acceptance is \num{0.259} of
cycles; and the misses are strongly clustered, with a mean blind run of
\num{9.16} cycles against the \num{3.53} that independent dropout would
give, and a longest blind run of \num{57} cycles, about \SI{16}{\second}.
That clustering is the point of the level. A planner that loses sight of
an obstacle for one cycle at a time behaves quite differently from one
that loses it for sixteen seconds, at identical average detection rates.

\subsection{S3: the vehicle, downstream of the planner}

S3 is a plant model applied to \texttt{/planner/cmd\_vel\_safe}, never to
the planner. It was identified from open-water trials of the physical
vehicle: full-command speeds of \SI{0.145}{\metre\per\second} in surge
and \SI{0.123}{\metre\per\second} in sway; yaw response of
\SI{0.0654}{\degree\per\second} per count with a 383-count deadband in
one direction and \num{0.0537} with 220 in the other; and first-order
rise and decay constants of \SI{2.0}{\second}/\SI{1.0}{\second} in surge
and \SI{0.30}{\second}/\SI{0.97}{\second} in sway.

Three honesty notes belong with those numbers. The translation axes are
reported as symmetric because the available trials cannot distinguish
their directions at this sample size, not because symmetry was verified;
\texttt{sway+} rests on a single valid trial. The rise constants have
spreads as large as their values. Command latency was never measured, so
the plant exposes a latency parameter and leaves it at zero rather than
injecting a plausible-looking value.

\subsection{What the ladder does not contain}

Two things measured during the campaign are deliberately absent. Pool
recirculation was measured at \SIrange{0.037}{0.053}{\metre\per\second}
and was non-stationary enough that successive measurements pointed in
opposite directions; fitting a disturbance model to four inconsistent
measurements would be invention, so the disturbance is an uncontrolled
input present in the physical runs and absent from every simulated one.
And, most consequentially, \emph{the ladder contains no visual detector}.
Simulated observations come from an oracle projection of the obstacle
geometry, degraded by the S1 and S2 statistics; every frozen run records
its perception source as an oracle relay and not a visual result. The
statistics were fitted to the real detector, so the surrogate is a
principled one, but it is a surrogate, and Section~\ref{sec:deployment}
reports what the real detector did that no statistical summary of it
predicted.
